\section{Discussion}
\label{sec:discussion}

\subsection{Key Findings}

\textbf{Finding 1: Efficiency without convergence.} The 10.2\% FLOP reduction derives from architectural iteration reduction, not attention convergence. Biological intuitions about refinement do not transfer directly.

\textbf{Finding 2: Iteration count dominates K/V caching.} The 10.2\% from $T$ reduction exceeds 6.8\% from caching.

\textbf{Finding 3: Constant scaling.} The 5.31\% reduction is sequence-length-independent.

\subsection{Limitations}

\begin{itemize}
    \item \textbf{L1:} Convergence tested on random-init model (needs trained weights)
    \item \textbf{L2:} CPU-only execution (no memory profiling)
    \item \textbf{L3:} Reasoning tasks (P3) not tested
    \item \textbf{L4:} Single architecture tested
\end{itemize}

\subsection{Broader Impact}

Our methodology promotes rigorous mechanism validation. Distinguishing ``what works'' from ``why it works'' enables principled architectural extensions. The methodology transfers to other efficiency claims.
